% ============================================================
% FICHIER DE CITATIONS — Mémoire TNAH
% Clotilde Long — 2026
%
%pour chaque citation, remplir les balises pour avoir un repère ------------------------------------------------------------
% Ce fichier n'est PAS importé dans main.tex
% C'est un carnet de travail personnel
%
% TYPES DE BALISES :
%   CITATION    → passage à citer quasi textuellement
%   PARAPHRASE  → idée à reformuler
%   DEF         → définition d'un concept clé
%   ARGUMENT    → soutient une de tes thèses
%   CONTRE      → position adverse, utile pour nuancer
%   CHIFFRE     → donnée chiffrée ou statistique
%
% EMPLACEMENTS PRÉVUS :
%   INTRO / ETAT-ART / METHODO / DISCUSSION / CONCLU
%
% RECHERCHE RAPIDE (Ctrl+F dans TeXstudio) :
%   → par thème    : ex. "ALGORITHME", "REQUETE", "CORPUS"
%   → par chapitre : ex. "ETAT-ART", "METHODO"
%   → par type     : ex. "DEF", "CITATION", "CONTRE"
% ============================================================


% %%%%%%%%%%%%%%%%%%%%%%%%%%%%%%%%%%%%%%%%%%%%%%%%%%%%%%%%%%%%
% THÈME 1 : INDEXATION AUTOMATIQUE
% %%%%%%%%%%%%%%%%%%%%%%%%%%%%%%%%%%%%%%%%%%%%%%%%%%%%%%%%%%%%


% ── DEF ─────────────────────────────────────────────────────
% TYPE        : DEF
% Thème       : INDEXATION / définition générale
% Source      : \footcite[p. ]{clé_zotero}
% Emplacement : ETAT-ART
% Statut      : @A-LIRE / @LU
%
% «Texte exact de la citation»
%
% → Reformulation / remarque personnelle :


% ── DEF ─────────────────────────────────────────────────────
% TYPE        : DEF
% Thème       : INDEXATION / indexation par mots-clés
% Source      : \footcite[p. ]{clé_zotero}
% Emplacement : ETAT-ART
% Statut      : @A-LIRE / @LU
%
% «Texte exact de la citation»
%
% → Reformulation / remarque personnelle :


% ── ARGUMENT ────────────────────────────────────────────────
% TYPE        : ARGUMENT
% Thème       : INDEXATION / limites de l'indexation manuelle
% Source      : \footcite[p. ]{clé_zotero}
% Emplacement : INTRO ou ETAT-ART
% Statut      : @A-LIRE / @LU
%
% «Texte exact de la citation»
%
% → Reformulation / remarque personnelle :


% ── CITATION ────────────────────────────────────────────────
% TYPE        : CITATION
% Thème       : INDEXATION / apprentissage automatique
% Source      : \footcite[p. ]{clé_zotero}
% Emplacement : ETAT-ART
% Statut      : @A-LIRE / @LU
%
% «Texte exact de la citation»
%
% → Reformulation / remarque personnelle :


% ── CHIFFRE ─────────────────────────────────────────────────
% TYPE        : CHIFFRE
% Thème       : INDEXATION / données chiffrées
% Source      : \footcite[p. ]{clé_zotero}
% Emplacement : INTRO ou DISCUSSION
% Statut      : @A-LIRE / @LU
%
% «Texte exact de la citation»
%
% → Reformulation / remarque personnelle :


% ── CONTRE ──────────────────────────────────────────────────
% TYPE        : CONTRE
% Thème       : INDEXATION / limites de l'indexation automatique
% Source      : \footcite[p. ]{clé_zotero}
% Emplacement : DISCUSSION
% Statut      : @A-LIRE / @LU
%
% «Texte exact de la citation»
%
% → Reformulation / remarque personnelle :
% → À relier avec :


% ── PARAPHRASE ──────────────────────────────────────────────
% TYPE        : PARAPHRASE
% Thème       : INDEXATION / évaluation des systèmes
% Source      : \footcite[p. ]{clé_zotero}
% Emplacement : METHODO
% Statut      : @A-LIRE / @LU
%
% «Texte exact de la citation»
%
% → Reformulation / remarque personnelle :


% ── ARGUMENT ────────────────────────────────────────────────
% TYPE        : ARGUMENT
% Thème       : INDEXATION / corpus et données d'entraînement
% Source      : \footcite[p. ]{clé_zotero}
% Emplacement : METHODO
% Statut      : @A-LIRE / @LU
%
% «Texte exact de la citation»
%
% → Reformulation / remarque personnelle :
% → À relier avec :




% %%%%%%%%%%%%%%%%%%%%%%%%%%%%%%%%%%%%%%%%%%%%%%%%%%%%%%%%%%%%
% THÈME 2 : RECHERCHE EN LANGAGE NATUREL
% %%%%%%%%%%%%%%%%%%%%%%%%%%%%%%%%%%%%%%%%%%%%%%%%%%%%%%%%%%%%


% ── DEF ─────────────────────────────────────────────────────
% TYPE        : DEF
% Thème       : LANGAGE-NATUREL / définition générale
% Source      : \footcite[p. ]{clé_zotero}
% Emplacement : ETAT-ART
% Statut      : @A-LIRE / @LU
%
% «Texte exact de la citation»
%
% → Reformulation / remarque personnelle :


% ── DEF ─────────────────────────────────────────────────────
% TYPE        : DEF
% Thème       : LANGAGE-NATUREL / traitement automatique du langage
% Source      : \footcite[p. ]{clé_zotero}
% Emplacement : ETAT-ART
% Statut      : @A-LIRE / @LU
%
% «Texte exact de la citation»
%
% → Reformulation / remarque personnelle :


% ── CITATION ────────────────────────────────────────────────
% TYPE        : CITATION
% Thème       : LANGAGE-NATUREL / requêtes en langue naturelle
% Source      : \footcite[p. ]{clé_zotero}
% Emplacement : ETAT-ART
% Statut      : @A-LIRE / @LU
%
% «Texte exact de la citation»
%
% → Reformulation / remarque personnelle :


% ── ARGUMENT ────────────────────────────────────────────────
% TYPE        : ARGUMENT
% Thème       : LANGAGE-NATUREL / avantages pour l'usager
% Source      : \footcite[p. ]{clé_zotero}
% Emplacement : INTRO ou ETAT-ART
% Statut      : @A-LIRE / @LU
%
% «Texte exact de la citation»
%
% → Reformulation / remarque personnelle :


% ── CONTRE ──────────────────────────────────────────────────
% TYPE        : CONTRE
% Thème       : LANGAGE-NATUREL / ambiguïté et limites
% Source      : \footcite[p. ]{clé_zotero}
% Emplacement : DISCUSSION
% Statut      : @A-LIRE / @LU
%
% «Texte exact de la citation»
%
% → Reformulation / remarque personnelle :
% → À relier avec :


% ── PARAPHRASE ──────────────────────────────────────────────
% TYPE        : PARAPHRASE
% Thème       : LANGAGE-NATUREL / modèles de langue
% Source      : \footcite[p. ]{clé_zotero}
% Emplacement : ETAT-ART ou METHODO
% Statut      : @A-LIRE / @LU
%
% «Texte exact de la citation»
%
% → Reformulation / remarque personnelle :


% ── CHIFFRE ─────────────────────────────────────────────────
% TYPE        : CHIFFRE
% Thème       : LANGAGE-NATUREL / performances des systèmes
% Source      : \footcite[p. ]{clé_zotero}
% Emplacement : DISCUSSION
% Statut      : @A-LIRE / @LU
%
% «Texte exact de la citation»
%
% → Reformulation / remarque personnelle :




% %%%%%%%%%%%%%%%%%%%%%%%%%%%%%%%%%%%%%%%%%%%%%%%%%%%%%%%%%%%%
% THÈME 3 : CROISEMENTS — INDEXATION + LANGAGE NATUREL
% %%%%%%%%%%%%%%%%%%%%%%%%%%%%%%%%%%%%%%%%%%%%%%%%%%%%%%%%%%%%

% (Citations qui articulent les deux thématiques de ton mémoire)


% ── ARGUMENT ────────────────────────────────────────────────
% TYPE        : ARGUMENT
% Thème       : CROISEMENT / langage naturel appliqué à l'indexation
% Source      : \footcite[p. ]{clé_zotero}
% Emplacement : INTRO ou CONCLU
% Statut      : @A-LIRE / @LU
%
% «Texte exact de la citation»
%
% → Reformulation / remarque personnelle :
% → Pourquoi c'est central pour ta problématique :


% ── CITATION ────────────────────────────────────────────────
% TYPE        : CITATION
% Thème       : CROISEMENT / systèmes hybrides
% Source      : \footcite[p. ]{clé_zotero}
% Emplacement : ETAT-ART ou DISCUSSION
% Statut      : @A-LIRE / @LU
%
% «Texte exact de la citation»
%
% → Reformulation / remarque personnelle :




% %%%%%%%%%%%%%%%%%%%%%%%%%%%%%%%%%%%%%%%%%%%%%%%%%%%%%%%%%%%%
% CITATIONS NON CLASSÉES — À TRIER
% %%%%%%%%%%%%%%%%%%%%%%%%%%%%%%%%%%%%%%%%%%%%%%%%%%%%%%%%%%%%

% (Coller ici les citations repérées rapidement,
%  à classer dans les thèmes plus tard)


% ── À CLASSER ───────────────────────────────────────────────
% TYPE        :
% Thème       :
% Source      : \footcite[p. ]{clé_zotero}
% Emplacement :
% Statut      : @A-LIRE / @LU
%
% «Texte exact de la citation»
%
% → Remarque :
